\chapter{Background and State of the Art}
\label{ch:background}

\section{Foundations of Machine Learning and Large Language Models}
\label{sec:ml-foundations}

\subsection{From classical statistical models to neural language models}
\label{sec:classical-to-neural}

A language model assigns a probability to a sequence of tokens, or equivalently predicts the next token given its predecessors. Early language models estimated this distribution with $n$-gram counts over a training corpus, a formulation that traces back to Shannon's information-theoretic treatment of language as a stochastic process \citep{shannon1948mathematical}. $n$-gram models are simple and exact but suffer from data sparsity: most sequences of length $n$ never occur in any finite corpus, forcing smoothing heuristics and capping the context that can be modeled.

Bengio et al. replaced discrete counts with a neural network that maps each vocabulary item to a dense, learned vector and predicts the next token from the concatenation of the preceding vectors \citep{bengio2003neural}. This established two ideas that persist through every subsequent architecture discussed in this chapter: (i) representing words as continuous embeddings rather than one-hot symbols, and (ii) training the representation and the predictive model jointly. Word2Vec \citep{mikolov2013distributed} and GloVe \citep{pennington2014glove} later factored the embedding step out as a standalone, context-independent pretraining objective, producing reusable word vectors later consumed by task-specific models. A limitation shared by all of these is that a word receives exactly one vector regardless of context — "bank" is encoded identically in "river bank" and "bank account." ELMo addressed this by deriving contextual embeddings from a bidirectional LSTM's hidden states, so a word's representation varies with its sentence \citep{peters2018deep}. This contextual property is a prerequisite for the Transformer architecture described next.

\subsection{Transformer architecture and autoregressive models}
\label{sec:transformer}

The Transformer replaced recurrence with self-attention: every token's representation is updated by attending to every other token in the sequence in parallel, weighted by learned compatibility scores \citep{vaswani2017attention}. Removing recurrence removes the sequential dependency that made recurrent networks slow to train on long sequences, and self-attention gives every pair of tokens a direct path to one another regardless of distance, which mitigates the vanishing-gradient behavior that limits how much context recurrent networks retain in practice.

Two families of models built on the Transformer diverge in how they use attention. Encoder models such as BERT attend bidirectionally — every token sees the whole input, including tokens to its right — and are trained with a masked-token prediction objective, making them suited to representation learning for classification and extraction tasks rather than open-ended generation \citep{devlin2019bert}. Decoder-only, autoregressive models such as the GPT family instead mask attention to the left context only, so each token is predicted from its predecessors alone, and generation proceeds one token at a time by repeatedly sampling from this conditional distribution \citep{radford2018improving,brown2020language}. This chapter uses "LLM" to refer to this decoder-only, autoregressive family, since it is the family the systems described in Chapter~\ref{ch:systems} consume through an API.

Scale is a defining property of this family. Training loss decreases as a smooth, predictable function of model size, dataset size, and compute, a relationship formalized as scaling laws \citep{kaplan2020scaling} and later refined by Chinchilla's finding that many early models were undertrained relative to their parameter count for the compute budget spent on them \citep{hoffmann2022training}. GPT-3 demonstrated that scale alone — without gradient updates at inference time — produces a model capable of performing new tasks from a handful of examples given in its input, a behavior termed few-shot in-context learning \citep{brown2020language}.

\subsection{General-purpose versus instruction-tuned LLMs; prompting}
\label{sec:instruction-tuning}

A base LLM trained only to predict the next token continues a prompt with what is statistically likely, which is not the same as what is requested. FLAN showed that fine-tuning on a large collection of tasks reformulated as natural-language instructions improves zero-shot performance on unseen tasks \citep{wei2022finetuned}. InstructGPT went further by fine-tuning first on human-written demonstrations and then on human preference comparisons via reinforcement learning, an alignment procedure known as RLHF that builds on preference-based reward learning \citep{christiano2017deep,ouyang2022training}. The resulting instruction-tuned models follow a natural-language instruction directly, which is the interaction mode every LLM call in the systems described in this thesis relies on: the caller writes a prompt, not a task-specific fine-tuning set.

Prompting is the practice of shaping a model's output through the text of the input alone, without updating any weights. Zero-shot prompting states the task with no examples; few-shot prompting prepends a small number of input-output demonstrations to the prompt \citep{brown2020language}. Chain-of-thought prompting elicits intermediate reasoning steps before the final answer, which measurably improves performance on multi-step reasoning tasks, and this improvement persists even with a single generic instruction ("let's think step by step") rather than task-specific demonstrations \citep{wei2022chain,kojima2022large}. Chain-of-thought is directly relevant to two design choices covered in Chapter~\ref{ch:systems}: self-consistency voting (\S\ref{sec:self-consistency}) operates on chain-of-thought samples, and the multi-hop retrieval loop in ColdRAG's inference engine scores and prunes reasoning paths in a structurally analogous way.

\subsection{Structured and constrained output from an LLM}
\label{sec:structured-output}

An autoregressive LLM's native output is unstructured text. Many downstream uses, including both systems in this thesis, require output that conforms to a fixed schema — a specific set of fields with specific types — so that it can be parsed and consumed programmatically without further natural-language interpretation. Two complementary strategies produce such output. The first constrains decoding itself: at each generation step, the set of tokens the model is allowed to sample is restricted to those consistent with a target grammar or JSON schema, guaranteeing well-formed output by construction rather than by post-hoc validation \citep{willard2023efficient}. The second strategy, function calling, exposes a schema to the model as part of the prompt and relies on instruction-following behavior (\S\ref{sec:instruction-tuning}) to produce a matching payload, with retries on validation failure rather than a hard decoding-time guarantee. Libraries such as Instructor layer this pattern on top of Pydantic, so a schema is defined once as a Pydantic model, passed to the LLM call, and the returned payload is parsed and validated against that same model, raising a validation error that can trigger an automatic retry with the error appended to the prompt \citep{liu2026instructor,colvin2026pydantic}. This is the mechanism Chapter~\ref{ch:systems} refers to concretely wherever a Pydantic model is generated at runtime and passed to an Instructor call. Structured-output reliability remains an active benchmarking concern rather than a solved problem, particularly for deeply nested or large schemas \citep{geng2025jsonschemabench}.

\subsection{Known limitations of LLMs relevant to this thesis}
\label{sec:llm-limitations}

Three limitations of LLMs motivate specific design choices in both systems described in this thesis and are referenced repeatedly in Chapter~\ref{ch:systems}.

\textbf{Hallucination.} An LLM can generate text that is fluent and confident but not supported by, or in direct conflict with, any grounding source \citep{ji2023survey,huang2023survey}. Detection approaches include comparing multiple sampled generations for mutual consistency \citep{manakul2023selfcheckgpt} and probing whether a model's own confidence correlates with correctness \citep{kadavath2022language,kuhn2023semantic}. The evidence-quote requirement described in Chapter~\ref{ch:systems} (every extracted field must carry a verbatim source quote, validated before acceptance) is a grounding mechanism aimed directly at this failure mode, applied at generation time rather than as post-hoc detection.

\textbf{Non-determinism and run-to-run variance.} Repeated calls to the same LLM with the same prompt and a temperature of zero are not guaranteed to return identical output: floating-point non-associativity under different batching and kernel-scheduling conditions on the serving side can change results between calls even when sampling itself is deterministic \citep{he2025defeating}. This matters for a system whose extraction pipeline is expected to be reproducible, and motivates measuring output variance directly (Chapter~\ref{ch:results}) rather than assuming determinism from a fixed temperature setting alone.

\textbf{The need for external grounding.} Because an LLM's parametric knowledge is fixed at training time and its generations are not intrinsically verifiable against a source, tasks that require up-to-date, document-specific, or verifiable facts cannot rely on the model's output alone. Retrieval-augmented generation, introduced next, is the dominant architectural response to this limitation.

\section{Retrieval-Augmented Generation and GraphRAG}
\label{sec:rag-graphrag}

\subsection{Retrieval-augmented generation: motivation, architecture, limits}
\label{sec:rag}

Retrieval-augmented generation (RAG) addresses the grounding limitation of \S\ref{sec:llm-limitations} by conditioning generation on text retrieved from an external corpus at inference time, rather than relying solely on parameters learned during training \citep{lewis2020retrieval}. The canonical architecture has two stages: a retriever selects a small set of passages relevant to the query, typically via dense vector similarity between a query embedding and a precomputed index of passage embeddings \citep{karpukhin2020dense}, and a generator conditions on the concatenation of the query and the retrieved passages to produce its answer. This decouples factual capacity from parameter count — the corpus, not the weights, is the source of truth — and lets the corpus be updated without retraining.

Flat RAG has two structural limits that motivate the graph-based extension in \S\ref{sec:graphrag}. First, retrieval treats passages as an unordered bag: similarity to the query is scored independently per passage, so nothing in the architecture represents relationships between passages, which makes questions that require combining evidence from several documents (multi-hop queries) hard to answer with a single retrieval pass \citep{yang2018hotpotqa}. Second, retrieval is local by construction — it surfaces passages similar to the query, not passages relevant to a broader thematic context the query does not literally mention, a limitation Lost in the Middle further compounds by showing that even correctly retrieved evidence is used unevenly by the generator depending on its position in the context window \citep{liu2024lost}. Multi-hop retrieval methods address the first limit by interleaving retrieval with intermediate reasoning steps, retrieving again based on what a partial chain of thought has inferred so far rather than only on the original query \citep{trivedi2023interleaving}.

\subsection{GraphRAG: knowledge graphs, community detection, and multi-hop reasoning}
\label{sec:graphrag}

GraphRAG replaces the flat passage index with a knowledge graph extracted from the corpus by an LLM, and replaces similarity-only retrieval with retrieval guided by graph structure \citep{edge2024local}. Two properties of a knowledge graph directly answer the two limits of flat RAG identified above: entities and relations are extracted once at ingestion time, which makes cross-document connections explicit as graph edges rather than implicit in embedding proximity, and traversal along edges gives retrieval a mechanism for multi-hop reasoning that a single similarity query cannot express.

Microsoft's GraphRAG builds this graph automatically — entities, relations, and claims are extracted from source documents by an LLM — and additionally applies the Leiden algorithm to partition the graph into hierarchical communities of densely connected nodes \citep{edge2024local,traag2019louvain}. Leiden improves on its predecessor Louvain by guaranteeing that every detected community is internally connected, which Louvain does not guarantee \citep{traag2019louvain}. Each community is then summarized by an LLM, and query answering proceeds by combining relevant community summaries — this is GraphRAG's "global search" mode, aimed at questions about the corpus as a whole rather than a specific fact. This community-summarization mechanism is not the retrieval pattern ColdRAG (\S\ref{sec:coldrag-paper}) or this thesis's systems use; the property carried forward from GraphRAG into this thesis's work is the more general one: a knowledge graph, built by an LLM from unstructured text, as the substrate that both grounds generation in explicit provenance (an answer can point to the specific graph edges that produced it) and supports iterative, multi-hop traversal rather than one-shot similarity retrieval. Subsequent work in this space — HippoRAG's memory-inspired indexing \citep{gutierrez2024hipporag} and Think-on-Graph's LLM-guided path exploration \citep{sun2024think} among others — explores different retrieval strategies over an LLM-built graph; \S\ref{sec:coldrag-paper} and \S\ref{sec:kg-construction} narrow this general family to the two specific frameworks — ColdRAG and DIAL-KG — this thesis's systems adapt directly.

\section{The Cold-Start Problem in Recommender Systems}
\label{sec:cold-start}

\subsection{Collaborative filtering, content-based, and hybrid approaches}
\label{sec:cf-content-hybrid}

Collaborative filtering (CF) recommends items to a user based on patterns in a user-item interaction matrix alone, under the assumption that users who agreed in the past will agree again: item-based CF scores a candidate item by its similarity, in terms of co-interaction, to items the user has already engaged with \citep{sarwar2001item}, and matrix factorization methods instead learn low-dimensional latent vectors for users and items such that their inner product approximates observed interaction strength \citep{koren2009matrix}. CF requires no information about item or user attributes, only interaction history, which is simultaneously its strength (domain-agnostic) and its structural weakness: an item or user with no interaction history has no signal for the model to factorize.

Content-based filtering recommends by comparing item or user attributes directly — describing an item by its features and recommending items similar in feature space to what a user has liked — which requires no interaction history and is therefore immune to the failure mode above, at the cost of needing informative, well-structured item attributes as an input. Hybrid systems combine the two, typically falling back to content-based signal when collaborative signal is absent or weak, and Burke's survey frames this combination along several axes (weighting, switching, feature combination) rather than treating it as a single fixed design \citep{burke2002hybrid}.

\subsection{Item cold-start versus user cold-start}
\label{sec:item-vs-user-coldstart}

Cold-start is the failure mode of collaborative filtering under sparse or absent interaction history, and it arises in two distinct forms. User cold-start occurs for a new user with no or minimal interaction history: there is nothing to factorize on the user side of the interaction matrix. Item cold-start is the symmetric case for a newly listed item with no interactions yet. The two are not interchangeable in either cause or remedy: user cold-start is commonly addressed by asking for explicit preferences or profile attributes at onboarding, whereas item cold-start (DropoutNet trains a model to reconstruct latent factors from content features alone, simulating the cold condition during training \citep{volkovs2017dropoutnet}; MeLU applies meta-learning so a model adapts to a new item or user from very few interactions \citep{lee2019melu}; Heater and CLCRec instead learn to align content-based and collaborative embedding spaces so a content-only representation is usable directly \citep{zhu2020recommendation,wei2021contrastive}) generally requires substituting content signal for interaction signal until enough interaction history accumulates. The system this thesis describes, ColdRAG, addresses item cold-start specifically: a newly ingested product has no interaction history by construction (\S\ref{sec:kg-construction}), and matching must proceed from its extracted content alone.

\subsection{LLM-based zero-shot recommenders}
\label{sec:llm-zeroshot-rec}

A separate line of work replaces learned collaborative or content embeddings with an LLM's parametric and in-context reasoning ability directly, framing recommendation as a language task. P5 unifies several recommendation subtasks (rating prediction, sequential recommendation, explanation) under one pretrained, prompted text-to-text model \citep{geng2022recommendation}; TALLRec fine-tunes an LLM on a small set of recommendation-specific instructions to align its general language ability with the recommendation task \citep{bao2023tallrec}; RecMind wraps an LLM as an agent that plans and calls external tools during recommendation, rather than answering from a single prompt \citep{wang2024recmind}. Independent of fine-tuning, LLMs prompted directly and without any task-specific training exhibit non-trivial zero-shot ranking ability \citep{dai2023uncovering,hou2024large}, and LLMRec augments this zero-shot capacity with a graph structure over items \citep{wei2024llmrec} — a combination structurally close to what \S\ref{sec:coldrag-paper} describes for ColdRAG. Recent surveys treat cold-start specifically as a setting where LLM-based, zero-training recommendation is particularly attractive, precisely because it requires no interaction-history-derived embedding for the cold item to begin with \citep{zhang2025cold,wu2024survey,lin2025how}.

\subsection{The ColdRAG paper and its relation to this thesis's work}
\label{sec:coldrag-paper}

Yang et al. propose ColdRAG, a retrieval-augmented framework for cold-start recommendation that builds a knowledge graph dynamically from item content and retrieves candidates through LLM-guided multi-hop reasoning over that graph rather than through a learned embedding trained on interaction history \citep{yang2025adaptive}. This combines the item-cold-start remedy of \S\ref{sec:item-vs-user-coldstart} (content substituting for absent interactions) with the GraphRAG retrieval pattern of \S\ref{sec:graphrag} (a knowledge graph as the retrieval substrate, traversed rather than searched flatly), applied specifically to the recommendation setting.

The system in this thesis, also named ColdRAG, is a domain-specific implementation adapted from this paper for real-estate off-market intermediation at Dometria AI; it is not the paper authors' original codebase, which is maintained separately. Chapter~\ref{ch:systems} describes this adaptation's engine (\S\ref{sec:coldrag-paper} above covers only the published framework this adaptation draws on) and states explicitly, at each point of divergence, where the Dometria implementation departs from the published design — most substantially in the ingestion pipeline, which follows DIAL-KG (\S\ref{sec:dial-kg}) rather than the construction method described in the ColdRAG paper itself.

\section{Knowledge Graph Construction}
\label{sec:kg-construction}

\subsection{Rule-based, supervised, and LLM-based approaches}
\label{sec:kgc-approaches}

Knowledge graph construction extracts structured entities and relations from unstructured text. Open information extraction systems such as TextRunner and ReVerb extract subject-relation-object triples using hand-built or lightly supervised syntactic patterns, with no fixed target schema \citep{banko2007open,fader2011identifying}. Distantly supervised approaches instead train a relation classifier by heuristically aligning an existing knowledge base with text — a sentence is labeled as expressing a relation if the entity pair it mentions is already related in that fashion in the knowledge base — trading manual annotation for a fixed, external schema \citep{mintz2009distant}.

LLM-based construction replaces both the hand-built patterns and the schema-aligned classifier with a prompted LLM that extracts entities and relations directly, using its pretrained language understanding rather than a task-specific supervised model \citep{zhu2024llms}. This shifts the construction bottleneck from labeled training data to prompt design and output validation, and introduces the reliability concerns of \S\ref{sec:llm-limitations} (hallucinated entities or relations, non-determinism across ingestion runs) directly into the graph-building step, which is the reason governance and evidence-checking mechanisms of the kind described in Chapter~\ref{ch:systems} (\S\ref{sec:kg-construction} continued below) sit downstream of extraction rather than being treated as optional.

\subsection{Schema-free and incremental knowledge graph construction}
\label{sec:schema-free-kgc}

A fixed schema, agreed before construction begins, is a strong assumption when the corpus is heterogeneous or evolving: a schema adequate for the first batch of documents may not accommodate entity types or relations that only appear later. Schema-free construction instead lets the schema emerge from the data, inferring entity and relation types as documents are processed rather than validating against a predefined ontology \citep{carta2023iterative}, and incremental construction processes documents one at a time or in small batches, merging each new extraction into a graph that already exists rather than building the graph once from a static corpus \citep{lairgi2024itext2kg}. The two properties compose naturally: an incremental pipeline needs a way to decide whether a newly extracted entity is the same as, or distinct from, an entity already in the graph, without a fixed schema to lean on for that decision — this is the entity disambiguation problem taken up in \S\ref{sec:ned}.

\subsection{The DIAL-KG framework}
\label{sec:dial-kg}

DIAL-KG, proposed by Bao et al., is a schema-free, incremental knowledge graph construction framework built around three components directly relevant to Chapter~\ref{ch:systems}'s description of ColdRAG's ingestion pipeline \citep{bao2026dial}. A Meta-Knowledge Base (MKB) accumulates entity profiles across ingestion runs, serving as the persistent context against which each new document's extraction is disambiguated and merged, rather than each document being processed in isolation. Dual-track extraction separates the extraction of graph-worthy entities and relations from the extraction of document- or product-level identifying information, so the two tracks can be validated and merged independently. An incremental adaptation cycle governs how the graph and the MKB evolve as new documents arrive: each ingestion both contributes new nodes and edges and updates the persistent entity profiles those nodes are checked against on the next ingestion.

This thesis's ColdRAG implementation adapts this framework as its ingestion pipeline (\S\ref{sec:kg-construction}, and in full in Chapter~\ref{ch:systems}); the MKB, dual-track extraction, and incremental cycle described here are the direct source of the corresponding named components in the implementation chapter — DIAL-KG is a framework this thesis's system implements against, not a component of ColdRAG's own published inference method (\S\ref{sec:coldrag-paper}).

\subsection{Named entity disambiguation, coreference resolution, entity linking}
\label{sec:ned}

Three related problems arise once entities are extracted and must be merged into a shared graph rather than kept per-document. Named entity disambiguation (NED) decides, given two entity mentions, whether they refer to the same real-world entity — the decision an incremental pipeline (\S\ref{sec:schema-free-kgc}) needs every time a new extraction is merged into an existing graph. Coreference resolution addresses the narrower, intra-document version of the same problem: deciding which mentions within a single text (including pronouns and other referring expressions) refer to the same entity \citep{lee2017end}. Entity linking (or entity resolution) is the version of the problem posed against an external, already-identified inventory of entities — is this mention entity 4711 in a known catalog? — commonly solved today via dense retrieval over entity representations followed by a re-ranking or classification step \citep{wu2020scalable}. All three tasks have moved from dedicated supervised models toward being posed as LLM prompts, though LLMs are not uniformly reliable coreference resolvers relative to specialized systems, particularly on cases requiring world knowledge outside the immediate text \citep{le2023large}, which motivates supplementing an LLM judge with a vector-similarity retrieval step rather than relying on the LLM's judgment alone \citep{sevgili2022neural} — the combination Chapter~\ref{ch:systems} describes as ColdRAG's NED mechanism.

\section{Self-Consistency and Validation via LLMs}
\label{sec:self-consistency-validation}

\subsection{Self-consistency and majority voting}
\label{sec:self-consistency}

Chain-of-thought prompting (\S\ref{sec:instruction-tuning}) produces one reasoning path per generation, and that single path can be wrong even when the model is, in aggregate, capable of reasoning correctly about the problem. Self-consistency addresses this by sampling multiple independent chain-of-thought completions for the same prompt at non-zero temperature and taking a majority vote over their final answers, discarding the individual reasoning paths and keeping only the answer each one converges to \citep{wang2023self}. Because errors in the reasoning process are not perfectly correlated across independent samples, but a correct final answer is a comparatively stable attractor that multiple different reasoning paths tend to reach, the majority vote is more reliable than any single sample. Wang et al. report consistent gains from applying self-consistency over standard chain-of-thought prompting across several reasoning benchmarks: GSM8K improves by 17.9 percentage points, SVAMP by 11.0, AQuA by 12.2, StrategyQA by 6.4, and ARC-challenge by 3.9 \citep{wang2023self}. Chapter~\ref{ch:systems} applies this same voting mechanism, restricted to mandatory and high-priority field groups rather than every extracted field, as a variance-reduction step in the extraction pipeline.

\subsection{LLM-as-judge and entailment gating}
\label{sec:llm-as-judge}

A separate use of an LLM is not to produce an answer but to evaluate one already produced — by itself, by another model, or by a preceding pipeline stage. LLM-as-judge substitutes an LLM's judgment for human evaluation or a fixed metric, prompting the model to score or compare candidate outputs against stated criteria \citep{zheng2023judging}. Judge reliability itself requires validation, since a judge is subject to the same limitations discussed in \S\ref{sec:llm-limitations}, and a substantial recent literature is specifically concerned with characterizing judge biases and agreement with human raters \citep{gu2024survey}.

Entailment gating is a specific, narrower application of the judge pattern to grounding: rather than scoring output quality along an open-ended rubric, the judge is asked a binary question — does this claimed value follow from (is it entailed by) a specific cited piece of source text? — and the claim is accepted only on a positive answer. This is a verification step applied after generation, distinct from constrained decoding (\S\ref{sec:structured-output}), which prevents ill-formed output at generation time but says nothing about whether well-formed output is actually supported by a source. Chapter~\ref{ch:systems} describes two concrete instances of this pattern: the governance layer's evidence check during knowledge graph construction (\S\ref{sec:kg-construction}), and the multi-hop reasoning engine's LLM-based edge scoring, which gates which graph edges a retrieval path is allowed to continue through.

\section{Vector Search and Embeddings}
\label{sec:vector-search}

\subsection{Embedding models and cosine similarity}
\label{sec:embedding-models}

A text embedding model maps a span of text to a fixed-length dense vector such that semantic similarity between two spans corresponds to geometric proximity between their vectors, most commonly measured by cosine similarity, the cosine of the angle between the two vectors, which is invariant to vector magnitude and therefore isolates directional (semantic) similarity from length effects. Modern embedding models are trained with a contrastive objective: pairs of texts known to be semantically related are pulled together in embedding space and unrelated pairs are pushed apart, a paradigm established for sentence-level embeddings by Sentence-BERT \citep{reimers2019sentence}. Multilingual embedding is a distinct extension of this problem, since a naively trained model has no reason to place a sentence and its translation nearby in embedding space; knowledge distillation from a monolingual teacher model into a multilingual student is one solution \citep{reimers2020making}, and BGE-M3 instead trains multilinguality, multiple retrieval granularities (sentence to long passage), and multiple retrieval functions (dense, sparse, multi-vector) jointly within a single model \citep{chen2024m3}. Chapter~\ref{ch:systems} names the two embedding models used by the systems in this thesis — a multilingual MiniLM distillation for contract-category matching, and BGE-M3 for ColdRAG's knowledge-graph entity embeddings — both instances of the families introduced here.

\subsection{Native vector indices and approximate nearest-neighbor search}
\label{sec:ann-indices}

Exact nearest-neighbor search over an embedding index requires comparing a query vector against every stored vector, which scales linearly with index size and becomes the retrieval bottleneck as a corpus grows. Approximate nearest-neighbor (ANN) methods trade a small, controllable amount of recall for large speedups. Hierarchical Navigable Small World graphs (HNSW) build a multi-layer proximity graph over the indexed vectors and search it greedily from a coarse top layer down to a fine bottom layer, achieving logarithmic-time search in practice at high recall \citep{malkov2020efficient}. Rather than a standalone ANN library, this thesis's ColdRAG system stores vectors as a native index inside its graph database, so retrieval and graph traversal (\S\ref{sec:graphrag}) operate over the same store without a separate synchronization step between a vector index and a graph store \citep{neo4j2026neo4j}.

\section{Human-in-the-Loop Data Acquisition}
\label{sec:hitl}

\subsection{Structured validation, approval gates, and confidence-based routing}
\label{sec:hitl-validation}

Human-in-the-loop (HITL) design keeps a person in the decision path at points where full automation is unreliable or unaccountable, rather than automating a task end-to-end \citep{amershi2014power}. Mixed-initiative interaction frames this as a shared control problem: the system and the user each contribute where they are more capable, and the interface must make it clear at each turn which party is expected to act \citep{horvitz1999principles}. Active learning is a related but distinct pattern, in which a system explicitly selects which items to route to a human (the most informative or most uncertain ones) rather than routing all or none \citep{settles2009active}; confidence-based routing in a conversational acquisition setting is the same selection principle applied per-field or per-turn, deciding whether a field is confident enough to accept automatically or must be confirmed with the user. In task-oriented conversational systems specifically, this often takes the form of a clarifying question issued when required information is missing or ambiguous rather than guessed \citep{aliannejadi2019asking,zamani2020generating} — the mechanism Chapter~\ref{ch:systems} describes for ColdRAG's conversational completion module, which asks a targeted follow-up question for each missing mandatory field rather than proceeding on an incomplete or inferred value.

\subsection{LlamaIndex Workflows: event-driven orchestration}
\label{sec:llamaindex-workflows}

LlamaIndex Workflows is an orchestration library structuring a multi-step process as a set of typed steps connected by typed events: each step declares which event type it consumes and which it emits, and the workflow engine dispatches control to whichever step's input type matches the event just produced, rather than following a single hand-coded control-flow graph \citep{liu2026llamaindex}. Events are themselves Pydantic models (\S\ref{sec:structured-output}), so the data passed between steps is schema-validated at each hand-off. A persistent Context object accompanies a workflow run and lets steps read and write shared state across turns, which is the property a multi-turn conversational acquisition process needs to accumulate partial answers across several user turns before all mandatory fields are filled. Chapter~\ref{ch:systems} describes ColdRAG's conversational completion module as an instance of exactly this pattern: a state machine over typed steps and events, with session state persisted between turns.

\section{Differential Privacy}
\label{sec:differential-privacy}

\subsection{$\varepsilon$-differential privacy, the Laplace mechanism, sensitivity, privacy budget}
\label{sec:dp-definition}

Differential privacy formalizes a guarantee about a data-releasing mechanism, not about a dataset: a mechanism is $\varepsilon$-differentially private if its output distribution changes by at most a factor of $e^{\varepsilon}$ whether or not any single individual's record is included in its input, so no output makes it possible to confidently infer whether a specific individual's data was used \citep{dwork2006differential,dwork2014algorithmic}. The parameter $\varepsilon$, the privacy budget, controls this trade-off directly: a smaller $\varepsilon$ bounds the distributional change more tightly and gives a stronger privacy guarantee at the cost of more distortion added to genuine query results, and a larger $\varepsilon$ does the reverse.

The Laplace mechanism achieves this guarantee for numeric queries by adding noise drawn from a Laplace distribution, calibrated to the query's sensitivity — the maximum amount the query's true output can change when one individual's record is added or removed — scaled by $1/\varepsilon$ \citep{dwork2006calibrating}. A query with low sensitivity (an individual record can barely move the result) tolerates less added noise for the same $\varepsilon$ than a high-sensitivity query, which is why sensitivity, not $\varepsilon$ alone, determines how much a specific numeric output is perturbed. Chapter~\ref{ch:systems} describes ColdRAG's application of exactly this mechanism to anonymize numeric metadata before it is exposed to a user during blind matching (\S\ref{sec:domain-context} below covers why blind matching is a domain requirement in the first place), with the privacy budget as a configurable parameter.

\section{Domain Context: Real-Estate Valuation and the Off-Market Sector}
\label{sec:domain-context}

\subsection{Valuation standards: RICS, IVS, EVS}
\label{sec:valuation-standards}

Real-estate valuation in Europe is governed by a small number of overlapping professional standards rather than a single statute. RICS's Red Book Global Standards set out valuation bases, mandatory reporting content, and professional conduct requirements for members of the Royal Institution of Chartered Surveyors \citep{royalinstitutionofcharteredsurveyors2024rics}. The International Valuation Standards (IVS), maintained by the International Valuation Standards Council, define valuation bases and methodology at a jurisdiction-agnostic level intended to be adopted or referenced by national standards \citep{internationalvaluationstandardscouncil2024international}. The European Valuation Standards (EVS), maintained by TEGoVA, occupy an intermediate position, harmonizing valuation practice specifically across European jurisdictions \citep{tegova2025european}. In Italy specifically, Tecnoborsa's Codice delle Valutazioni Immobiliari provides the domain-specific standard that a document-extraction pipeline operating over Italian real-estate documents (Chapter~\ref{ch:systems}'s bulk-insert system) must be consistent with in its choice of extracted fields and valuation terminology \citep{tecnoborsa2025codice}. This thesis's systems do not implement valuation calculations against any of these standards; they are cited here because the extraction schema and terminology used by the bulk-insert system (Chapter~\ref{ch:systems}) is designed to align with the field and category conventions these standards establish, not because the systems perform standard-compliant valuation themselves.

\subsection{Blind matching, confidential deal-flow, VDRs, and NDAs}
\label{sec:blind-matching}

Off-market intermediation — matching a real-estate asset with a counterparty before or without a public listing — operates under confidentiality constraints that public-market matching does not. A seller's identity, an asset's precise location, or the fact that an asset is for sale at all is frequently sensitive information, disclosed to a prospective counterparty only incrementally as a deal progresses, commonly gated by a non-disclosure agreement (NDA) and, once due diligence begins, exchanged through a virtual data room (VDR) rather than openly. "Blind" matching is the practice of presenting a prospective match to a counterparty using only non-identifying summary attributes, deferring full disclosure until both sides have expressed mutual interest.

This confidentiality requirement is the direct, practical motivation for the differential-privacy anonymization step described in \S\ref{sec:dp-definition} and, in implementation, in Chapter~\ref{ch:systems}: a matching system that streams candidate metadata to a user during exploration must not let numeric metadata (values that could, in aggregate across many queries, re-identify a specific listing) leak more information than the blind-matching convention permits. This confidentiality problem is a specific instance of a more general information-asymmetry condition in markets with hidden characteristics \citep{akerlof1970market}, and the underlying question a matching system answers — pairing two sides of a market on the basis of mutual preference — has its formal roots in two-sided matching theory \citep{gale1962college,roth1990two}.
